\chapter{Boolean Analysis — Bool BLR}
%%% generated by the blueprint pipeline from TCSlib.BooleanAnalysis.BLR.BoolBLR
\section{Overview}

This module develops the Blum--Luby--Rubinfeld (BLR) linearity test for
Boolean functions $f : \{0,1\}^n \to \{0,1\}$.  It establishes that a
function is linear (i.e.\ satisfies $f(x \oplus y) = f(x) \oplus f(y)$)
if and only if its $\pm 1$ lift equals some Fourier character $\chi_S$, and
proves both completeness (a linear function passes BLR with probability $1$)
and soundness (a function $\varepsilon$-far from every linear function passes
BLR with probability at most $1-\varepsilon$), with the soundness bound
derived via a Fourier-analytic argument using Parseval's identity.

%%%%%%%%%%%%%%%%%%%%%%%%%%%%%%%%%%%%%%%%%%%%%%%%%%%%%%%%%%%%%%%%%%%%%%%%%%%%%%%
\section{Declarations}

\begin{definition}[Linearity of a Boolean function]
\lean{BoolBLR.is_linear_bool}
\label{BoolBLR.is_linear_bool}
\leanok
\uses{BoolFourier.hypercube, BoolFourier.xor_vec}
A function $f : \{0,1\}^n \to \{0,1\}$ is \emph{linear} if for all
$x, y \in \{0,1\}^n$ we have
\[
  f(x \oplus y) = f(x) \oplus f(y),
\]
where $\oplus$ denotes bitwise XOR (i.e.\ \texttt{Bool.xor}).
\end{definition}

\begin{definition}[$\pm 1$ lift of a Boolean function]
\lean{BoolBLR.lift_pm1}
\label{BoolBLR.lift_pm1}
\leanok
\uses{BoolFourier.BoolFun, BoolFourier.BoolToPM1, BoolFourier.hypercube}
Given $f : \{0,1\}^n \to \{0,1\}$, its \emph{$\pm 1$ lift}
$\texttt{BoolBLR.lift\_pm1}\,f : \{0,1\}^n \to \mathbb{R}$ is defined by
$x \mapsto (-1)^{f(x)}$, converting each Boolean output to a real sign via
\texttt{BoolToPM1}.
\end{definition}

\begin{definition}[Boolean distance]
\lean{BoolBLR.bool_dist}
\label{BoolBLR.bool_dist}
\leanok
\uses{BoolFourier.expectation, BoolFourier.hypercube}
The \emph{distance} between two Boolean functions
$f, g : \{0,1\}^n \to \{0,1\}$ is
\[
  \dist(f, g) \;=\; \Pr_{x \sim \{0,1\}^n}[f(x) \ne g(x)],
\]
computed as the uniform expectation of the indicator $\mathbf{1}[f(x) \ne g(x)]$.
\end{definition}

\begin{definition}[$\varepsilon$-far from linear]
\lean{BoolBLR.epsilon_far_from_linear}
\label{BoolBLR.epsilon_far_from_linear}
\leanok
\uses{BoolBLR.bool_dist, BoolBLR.is_linear_bool, BoolFourier.hypercube}
A function $f : \{0,1\}^n \to \{0,1\}$ is \emph{$\varepsilon$-far from linear}
if $0 \le \varepsilon \le 1$ and for every linear function
$g : \{0,1\}^n \to \{0,1\}$ we have $\dist(f, g) \ge \varepsilon$.
\end{definition}

\begin{lemma}[Linear Boolean functions are exactly the Walsh characters]
\lean{BoolBLR.linear_bool_iff_character}
\label{BoolBLR.linear_bool_iff_character}
\leanok
\uses{BoolBLR.is_linear_bool, BoolBLR.lift_pm1, BoolFourier.BoolToPM1, BoolFourier.BoolToPM1_xor, BoolFourier.char_S, BoolFourier.hypercube, BoolFourier.xor_vec, BooleanAnalysis.boolToSign, BooleanAnalysis.chiS}
\difficulty{6}
Let $f : \{0,1\}^n \to \{0,1\}$ be a Boolean function. Then $f$ is linear — that is,
$f(x \oplus y) = f(x) \oplus f(y)$ for all $x, y \in \{0,1\}^n$, where $\oplus$ denotes
bitwise XOR — if and only if there exists a set $S \subseteq [n]$ for which the $\pm 1$
lift $x \mapsto (-1)^{f(x)}$ coincides with the Walsh character $\chi_S(x) = \prod_{i
\in S}(-1)^{x_i}$.
\end{lemma}

\begin{definition}[BLR acceptance probability]
\lean{BoolBLR.BLR_accept_prob}
\label{BoolBLR.BLR_accept_prob}
\leanok
\uses{BoolBLR.lift_pm1, BoolFourier.expectation, BoolFourier.hypercube, BoolFourier.xor_vec}
The \emph{BLR acceptance probability} of $f : \{0,1\}^n \to \{0,1\}$ is
\[
  \Pr_{x,y}[\text{BLR accepts}]
  \;=\;
  \E_{x}\!\left[\E_{y}\!\left[
    \mathbf{1}\!\left[
      \texttt{BoolBLR.lift\_pm1}\,f(x \oplus y)
      = \texttt{BoolBLR.lift\_pm1}\,f(x)\cdot\texttt{BoolBLR.lift\_pm1}\,f(y)
    \right]
  \right]\right],
\]
where $x, y$ are drawn uniformly from $\{0,1\}^n$.
\end{definition}

\begin{lemma}[BLR acceptance probability in ±1 form]
\lean{BoolBLR.BLR_accept_prob_pm1}
\label{BoolBLR.BLR_accept_prob_pm1}
\leanok
\uses{BoolBLR.BLR_accept_prob, BoolBLR.lift_pm1, BoolFourier.card_hypercube, BoolFourier.expectation, BoolFourier.hypercube, BoolFourier.xor_vec}
\difficulty{3}
Let $f : \{0,1\}^n \to \{0,1\}$ be a Boolean-valued function, and write $F(x) =
(-1)^{f(x)} \in \{\pm 1\}$ for its sign lift. Then the BLR acceptance probability of
$f$, the probability that $F(x \oplus y) = F(x)\,F(y)$ when $x$ and $y$ are drawn
independently and uniformly from $\{0,1\}^n$, satisfies
\[
  \Pr_{x,y}[\text{BLR accepts}]
  \;=\;
  \frac{1 + \E_{x}\!\left[\E_{y}\!\left[ F(x)\,F(y)\,F(x \oplus y) \right]\right]}{2},
\]
where $x \oplus y$ denotes the componentwise XOR.
\end{lemma}

\begin{lemma}[Triple expectation as an inner expectation]
\lean{BoolBLR.triple_expectation_as_convolution}
\label{BoolBLR.triple_expectation_as_convolution}
\leanok
\uses{BoolFourier.BoolFun, BoolFourier.expectation, BoolFourier.xor_vec}
\difficulty{2}
Let $f : \{0,1\}^n \to \bbr$ be a Boolean function on the hypercube, write $x \oplus y$
for the componentwise XOR of $x, y \in \{0,1\}^n$, and let $\E$ denote expectation over
a uniformly random point of $\{0,1\}^n$. Then
\[
  \E_x\!\left[\E_y\!\left[f(x)\,f(y)\,f(x \oplus y)\right]\right]
  \;=\;
  \E_x\!\left[f(x)\cdot\E_y\!\left[f(y)\,f(x \oplus y)\right]\right].
\]
\end{lemma}

\begin{lemma}[Cube of Fourier coefficients as an autocorrelation expectation]
\lean{BoolBLR.triple_expectation_eq_cube_fourier}
\label{BoolBLR.triple_expectation_eq_cube_fourier}
\leanok
\uses{BoolBLR.lift_pm1, BoolFourier.char_S, BoolFourier.convolution, BoolFourier.expectation, BoolFourier.fourier_coeff, BoolFourier.fourier_coeff_convolution, BoolFourier.fourier_expansion, BoolFourier.hypercube, BoolFourier.xor_vec}
\difficulty{3}
Let $f : \{0,1\}^n \to \{0,1\}$ be a Boolean-valued function on the hypercube, and let
$g = (-1)^{f}$ denote its $\pm 1$ lift. Then the double uniform expectation of the
triple product of $g$ over independent inputs and their componentwise XOR equals the sum
of the cubes of the Fourier coefficients of $g$:
\[
  \E_{x}\!\left[\,\E_{y}\!\left[\,g(x)\,g(y)\,g(x \oplus y)\,\right]\right]
  \;=\;
  \sum_{S \subseteq [n]} \hat{g}(S)^{3},
\]
where $x, y$ range uniformly over $\{0,1\}^n$, $(x \oplus y)_i = x_i \oplus y_i$, and
$\hat{g}(S) = \langle g, \chi_S \rangle$ is the Fourier coefficient at $S \subseteq
[n]$.
\end{lemma}

\begin{lemma}[Fourier coefficient bound from distance]
\lean{BoolBLR.fourier_coeff_le_of_dist_ge}
\label{BoolBLR.fourier_coeff_le_of_dist_ge}
\leanok
\uses{BoolBLR.bool_dist, BoolBLR.lift_pm1, BoolFourier.BoolToPM1, BoolFourier.card_hypercube, BoolFourier.char_S, BoolFourier.expectation, BoolFourier.fourier_coeff, BoolFourier.hypercube, BoolFourier.inner_product}
\difficulty{4}
Let $f,g:\{0,1\}^n\to\{0,1\}$ be Boolean-valued functions, let $S\subseteq[n]$, and let
$\varepsilon\in\bbr$. Write $f^{\pm},g^{\pm}:\{0,1\}^n\to\bbr$ for their $\pm1$ lifts,
given by $x\mapsto(-1)^{f(x)}$ and $x\mapsto(-1)^{g(x)}$. Suppose that $g^{\pm}$
coincides with the Walsh character $\chi_S$, and that $\dist(f,g)=\Pr_{x}[f(x)\ne
g(x)]\ge\varepsilon$. Then the Fourier coefficient of $f^{\pm}$ at $S$ satisfies
\[\widehat{f^{\pm}}(S)\;\le\;1-2\varepsilon.\]
\end{lemma}

\begin{lemma}[Fourier coefficient bound for functions far from linear]
\lean{BoolBLR.fourier_coeff_le_of_far_from_linear}
\label{BoolBLR.fourier_coeff_le_of_far_from_linear}
\leanok
\uses{BoolBLR.epsilon_far_from_linear, BoolBLR.fourier_coeff_le_of_dist_ge, BoolBLR.is_linear_bool, BoolBLR.lift_pm1, BoolBLR.linear_bool_iff_character, BoolFourier.char_S, BoolFourier.fourier_coeff, BoolFourier.hypercube, BooleanAnalysis.chiS}
\difficulty{5}
Let $f : \{0,1\}^n \to \{0,1\}$ be $\varepsilon$-far from linear, meaning that $0 \le
\varepsilon \le 1$ and every linear function $g : \{0,1\}^n \to \{0,1\}$ (one satisfying
$g(x \oplus y) = g(x) \oplus g(y)$ for all $x,y$) has $\Pr_{x}[f(x) \ne g(x)] \ge
\varepsilon$ under the uniform distribution. Let $F(x) = (-1)^{f(x)}$ be the $\pm 1$
lift of $f$. Then for every subset $S \subseteq [n]$, the Fourier coefficient of $F$ at
$S$ satisfies
\[
  \hat F(S) \;=\; \E_{x}\bigl[F(x)\,\chi_S(x)\bigr] \;\le\; 1 - 2\varepsilon,
\]
where $\chi_S(x) = \prod_{i \in S}(-1)^{x_i}$ is the Walsh character indexed by $S$.
\end{lemma}

\begin{lemma}[Sum of cubed Fourier coefficients bounds distance to linearity]
\lean{BoolBLR.BLR_soundness_via_fourier}
\label{BoolBLR.BLR_soundness_via_fourier}
\leanok
\uses{BoolBLR.epsilon_far_from_linear, BoolBLR.fourier_coeff_le_of_far_from_linear, BoolBLR.lift_pm1, BoolFourier.BoolToPM1_sq, BoolFourier.L2_norm_sq, BoolFourier.expectation, BoolFourier.fourier_coeff, BoolFourier.hypercube, BoolFourier.parseval_identity}
\difficulty{3}
Let $f : \{0,1\}^n \to \{0,1\}$ be a Boolean function and let $\varepsilon \in \bbr$
with $0 \le \varepsilon \le 1$, and suppose $f$ is $\varepsilon$-far from linear: for
every linear function $g : \{0,1\}^n \to \{0,1\}$ one has $\dist(f,g) \ge \varepsilon$.
Writing $F = (-1)^{f}$ for the $\pm 1$ lift of $f$ and $\hat F(S) = \langle F,
\chi_S\rangle$ for its Fourier coefficient at $S \subseteq [n]$, the sum of the cubed
coefficients over all subsets satisfies
\[
  \sum_{S \subseteq [n]} \hat F(S)^3 \;\le\; 1 - 2\varepsilon.
\]
\end{lemma}

\begin{lemma}[Completeness of the BLR linearity test]
\lean{BoolBLR.BLR_completeness}
\label{BoolBLR.BLR_completeness}
\leanok
\uses{BoolBLR.BLR_accept_prob, BoolBLR.is_linear_bool, BoolBLR.lift_pm1, BoolFourier.BoolToPM1_xor, BoolFourier.expectation, BoolFourier.hypercube, BoolFourier.xor_vec}
\difficulty{2}
Let $n$ be a natural number and let $f : \{0,1\}^n \to \{0,1\}$ be linear, meaning that
$f(x \oplus y) = f(x) \oplus f(y)$ for all $x, y \in \{0,1\}^n$. Then the BLR acceptance
probability of $f$ equals $1$; that is, when $x$ and $y$ are drawn uniformly and
independently from $\{0,1\}^n$, the test always accepts.
\end{lemma}

\begin{lemma}[Soundness of the BLR linearity test]
\lean{BoolBLR.BLR_soundness}
\label{BoolBLR.BLR_soundness}
\leanok
\uses{BoolBLR.BLR_accept_prob, BoolBLR.BLR_accept_prob_pm1, BoolBLR.BLR_soundness_via_fourier, BoolBLR.epsilon_far_from_linear, BoolBLR.triple_expectation_eq_cube_fourier, BoolFourier.hypercube}
\difficulty{2}
Let $f : \{0,1\}^n \to \{0,1\}$ and let $\varepsilon \in \bbr$ satisfy $0 \le
\varepsilon \le 1$. Suppose $f$ is $\varepsilon$-far from linear, meaning that every
linear function $g : \{0,1\}^n \to \{0,1\}$ (one satisfying $g(x \oplus y) = g(x) \oplus
g(y)$ for all $x, y$) has distance $\Pr_{x}[f(x) \ne g(x)] \ge \varepsilon$ from $f$.
Then the BLR acceptance probability of $f$ satisfies
\[
\E_{x}\!\left[\E_{y}\!\left[\mathbf{1}\!\left[(-1)^{f(x \oplus y)} =
(-1)^{f(x)}\,(-1)^{f(y)}\right]\right]\right] \;\le\; 1 - \varepsilon,
\]
where $x$ and $y$ range uniformly over $\{0,1\}^n$; equivalently, the test rejects $f$
with probability at least $\varepsilon$.
\end{lemma}

%%%%%%%%%%%%%%%%%%%%%%%%%%%%%%%%%%%%%%%%%%%%%%%%%%%%%%%%%%%%%%%%%%%%%%%%%%%%%%%
%%% added by the blueprint pipeline
\section{Additional declarations}

\begin{lemma}[Value of a linear Boolean function on an indicator vector]
\lean{BoolBLR.linear_bool_iff_character_aux_h_fx_1}
\label{BoolBLR.linear_bool_iff_character_aux_h_fx_1}
\leanok
\uses{BoolBLR.is_linear_bool, BoolFourier.hypercube, BoolFourier.xor_vec}
\difficulty{5}
Let $f : \{0,1\}^n \to \{0,1\}$ be a linear Boolean function, meaning $f(x \oplus y) =
f(x) \oplus f(y)$ for all $x, y \in \{0,1\}^n$, where $\oplus$ is componentwise XOR. For
a subset $s \subseteq \{0, \dots, n-1\}$, let $\mathbf{1}_s \in \{0,1\}^n$ be its
indicator vector, whose $i$-th coordinate is $1$ when $i \in s$ and $0$ otherwise, and
for each index $i$ let $e_i \in \{0,1\}^n$ be the standard basis vector, whose $i$-th
coordinate is $1$ and all others $0$. Then for every $s$ the value $f(\mathbf{1}_s)$ is
$0$ when
\[
  \sum_{i \in s} \mathbf{1}[\,f(e_i) = 1\,] \equiv 0 \pmod 2,
\]
and $1$ otherwise; that is, $f(\mathbf{1}_s)$ equals the parity of the number of $i \in
s$ with $f(e_i) = 1$.
\end{lemma}

\begin{lemma}[Value of a linear Boolean function from its values on basis vectors]
\lean{BoolBLR.linear_bool_iff_character_aux_h_fx_2}
\label{BoolBLR.linear_bool_iff_character_aux_h_fx_2}
\leanok
\uses{BoolBLR.is_linear_bool, BoolBLR.linear_bool_iff_character_aux_h_fx_1, BoolFourier.hypercube, BooleanAnalysis.BoolCube}
\difficulty{2}
Let $f : \{0,1\}^n \to \{0,1\}$ be a linear Boolean function, meaning $f(x \oplus y) =
f(x) \oplus f(y)$ for all $x, y \in \{0,1\}^n$, and for each coordinate $i$ let $e_i \in
\{0,1\}^n$ denote the standard basis vector that is $1$ in position $i$ and $0$
elsewhere. Then for every $x \in \{0,1\}^n$, the value $f(x)$ equals the parity of the
number of coordinates $i$ with $x_i = 1$ and $f(e_i) = 1$; that is, $f(x) = 1$ if that
count is odd and $f(x) = 0$ if it is even.
\end{lemma}

\begin{lemma}[Multiplicativity of the Walsh characters]
\lean{BoolBLR.linear_bool_iff_character_aux_h_char}
\label{BoolBLR.linear_bool_iff_character_aux_h_char}
\leanok
\uses{BoolFourier.BoolToPM1_xor, BoolFourier.char_S, BoolFourier.hypercube, BoolFourier.xor_vec, BooleanAnalysis.chiS}
\difficulty{3}
Fix $n$ and a subset $S \subseteq [n]$, and let $\chi_S(x) = \prod_{i \in S}(-1)^{x_i}$
be the associated Walsh character on the hypercube $\{0,1\}^n$. Then for all $x, y \in
\{0,1\}^n$,
\[
  \chi_S(x \oplus y) \;=\; \chi_S(x)\,\chi_S(y),
\]
where $x \oplus y$ denotes the componentwise XOR.
\end{lemma}

\begin{lemma}[XOR-additivity of a Boolean function whose lift is a character]
\lean{BoolBLR.linear_bool_iff_character_aux_h_eq}
\label{BoolBLR.linear_bool_iff_character_aux_h_eq}
\leanok
\uses{BoolBLR.lift_pm1, BoolFourier.BoolToPM1, BoolFourier.BoolToPM1_xor, BoolFourier.char_S, BoolFourier.hypercube, BoolFourier.xor_vec}
\difficulty{2}
Let $f : \{0,1\}^n \to \{0,1\}$ be a Boolean-valued function whose $\pm 1$ lift
coincides with a Walsh character, say $x \mapsto (-1)^{f(x)}$ equals $\chi_S$ for some
$S \subseteq [n]$, and fix points $x, y \in \{0,1\}^n$. If the character is
multiplicative at this pair, meaning $\chi_S(x \oplus y) = \chi_S(x)\,\chi_S(y)$ where
$x \oplus y$ denotes the componentwise XOR, then
\[
  (-1)^{f(x \oplus y)} \;=\; (-1)^{f(x) \oplus f(y)}.
\]
\end{lemma}

\begin{lemma}[Value of a linear Boolean function from its action on basis vectors]
\lean{BoolBLR.linear_bool_iff_character_aux_h_fx}
\label{BoolBLR.linear_bool_iff_character_aux_h_fx}
\leanok
\uses{BoolBLR.is_linear_bool, BoolBLR.linear_bool_iff_character_aux_h_fx_2, BoolFourier.hypercube, BooleanAnalysis.BoolCube}
\difficulty{1}
Let $f : \{0,1\}^n \to \{0,1\}$ be a linear Boolean function, meaning $f(x \oplus y) =
f(x) \oplus f(y)$ for all $x, y \in \{0,1\}^n$, where $\oplus$ is componentwise XOR. For
each coordinate $i$, let $e_i \in \{0,1\}^n$ denote the standard basis vector that is
$1$ in coordinate $i$ and $0$ elsewhere. Then for every $x \in \{0,1\}^n$, the value
$f(x)$ is $1$ if the number of coordinates $i$ with $x_i = 1$ and $f(e_i) = 1$ is odd,
and $0$ if that number is even; equivalently, $f(x)$ is the parity of $\lvert\{\, i :
x_i = 1 \text{ and } f(e_i) = 1 \,\}\rvert$.
\end{lemma}

\begin{lemma}[XOR-additivity of a Boolean function that lifts to a character]
\lean{BoolBLR.linear_bool_iff_character_aux_h_eq_h}
\label{BoolBLR.linear_bool_iff_character_aux_h_eq_h}
\leanok
\uses{BoolBLR.lift_pm1, BoolBLR.linear_bool_iff_character_aux_h_char, BoolBLR.linear_bool_iff_character_aux_h_eq, BoolFourier.BoolToPM1, BoolFourier.char_S, BoolFourier.hypercube, BoolFourier.xor_vec}
\difficulty{1}
Let $f : \{0,1\}^n \to \{0,1\}$ be a Boolean-valued function whose $\pm 1$ lift $x
\mapsto (-1)^{f(x)}$ coincides with the Walsh character $\chi_S$ for some subset $S
\subseteq [n]$. Then for all $x, y \in \{0,1\}^n$,
\[
  (-1)^{f(x \oplus y)} \;=\; (-1)^{f(x) \oplus f(y)},
\]
where $x \oplus y$ denotes the componentwise XOR of $x$ and $y$ and $f(x) \oplus f(y)$
the XOR of the two Boolean outputs.
\end{lemma}

\begin{lemma}[Pointwise $\pm 1$ form of the BLR indicator]
\lean{BoolBLR.BLR_accept_prob_pm1_aux_h_eq}
\label{BoolBLR.BLR_accept_prob_pm1_aux_h_eq}
\leanok
\uses{BoolBLR.lift_pm1, BoolFourier.hypercube, BoolFourier.xor_vec}
\difficulty{3}
Fix $n \in \bbn$ and a Boolean-valued function $f : \{0,1\}^n \to \{0,1\}$, and let $g :
\{0,1\}^n \to \bbr$ be its $\pm 1$ lift, $g(x) = (-1)^{f(x)}$. Then for all $x, y \in
\{0,1\}^n$,
\[
  \mathbf{1}\!\left[\, g(x \oplus y) = g(x)\,g(y) \,\right]
  \;=\;
  \frac{1 + g(x)\,g(y)\,g(x \oplus y)}{2},
\]
where $x \oplus y$ denotes the componentwise XOR and the left-hand side is $1$ when the
equality holds and $0$ otherwise.
\end{lemma}

\begin{lemma}[Fourier expansion of the self-convolution]
\lean{BoolBLR.triple_expectation_eq_cube_fourier_aux_h_convolution}
\label{BoolBLR.triple_expectation_eq_cube_fourier_aux_h_convolution}
\leanok
\uses{BoolBLR.lift_pm1, BoolFourier.char_S, BoolFourier.convolution, BoolFourier.fourier_coeff, BoolFourier.fourier_coeff_convolution, BoolFourier.fourier_expansion, BoolFourier.hypercube, BooleanAnalysis.BoolCube}
\difficulty{3}
Let $f : \{0,1\}^n \to \{0,1\}$ be a Boolean-valued function, and let $g : \{0,1\}^n \to
\bbr$ be its $\pm 1$ lift, $g(x) = (-1)^{f(x)}$. Writing $g * g$ for the
self-convolution $(g * g)(x) = \E_{y}\bigl[g(y)\,g(x \oplus y)\bigr]$ with respect to
the uniform measure, and $\hat g(S) = \E[g\,\chi_S]$ for the Fourier coefficient of $g$
at $S$, one has for every $x \in \{0,1\}^n$
\[
  (g * g)(x) \;=\; \sum_{S \subseteq [n]} \hat g(S)^2\, \chi_S(x),
\]
where $\chi_S(x) = \prod_{i \in S} (-1)^{x_i}$ is the Walsh character of $S$. That is,
the self-convolution of $g$ has Fourier coefficients $\hat g(S)^2$.
\end{lemma}

\begin{lemma}[Expectation of a function against its Fourier-expanded autocorrelation]
\lean{BoolBLR.triple_expectation_eq_cube_fourier_aux_h_substitute}
\label{BoolBLR.triple_expectation_eq_cube_fourier_aux_h_substitute}
\leanok
\uses{BoolBLR.lift_pm1, BoolFourier.char_S, BoolFourier.convolution, BoolFourier.expectation, BoolFourier.fourier_coeff, BoolFourier.hypercube, BooleanAnalysis.BoolCube}
\difficulty{3}
Let $f : \{0,1\}^n \to \{0,1\}$ be a Boolean function, and let $g$ denote its $\pm 1$
lift, so that $g(x) = (-1)^{f(x)}$. Suppose that the convolution of $g$ with itself
admits the Fourier expansion
\[
  (g * g)(x) \;=\; \sum_{S \subseteq [n]} \hat{g}(S)^2\,\chi_S(x)
  \qquad\text{for every } x \in \{0,1\}^n,
\]
where $\chi_S$ is the Walsh character of $S$ and $\hat{g}(S) = \langle g,
\chi_S\rangle$. Then
\[
  \E_x\!\left[g(x)\,(g * g)(x)\right]
  \;=\;
  \sum_{S \subseteq [n]} \hat{g}(S)^2\,\E_x\!\left[g(x)\,\chi_S(x)\right],
\]
the expectation being taken uniformly over $x \in \{0,1\}^n$.
\end{lemma}

\begin{lemma}[Product of $\pm 1$ lifts as a disagreement indicator]
\lean{BoolBLR.fourier_coeff_le_of_dist_ge_aux_h_lift_pm1}
\label{BoolBLR.fourier_coeff_le_of_dist_ge_aux_h_lift_pm1}
\leanok
\uses{BoolBLR.lift_pm1, BoolFourier.BoolToPM1, BoolFourier.hypercube, BooleanAnalysis.BoolCube}
\difficulty{2}
Let $f, g : \{0,1\}^n \to \{0,1\}$ be Boolean-valued functions on the hypercube, and
write $F, G : \{0,1\}^n \to \bbr$ for their $\pm 1$ lifts, given by $F(x) = (-1)^{f(x)}$
and $G(x) = (-1)^{g(x)}$. Then for every point $x \in \{0,1\}^n$,
\[
  F(x)\,G(x) \;=\; 1 - 2\cdot\mathbf{1}\!\left[f(x) \ne g(x)\right],
\]
so the product equals $+1$ where $f$ and $g$ agree at $x$ and $-1$ where they disagree.
\end{lemma}

\begin{lemma}[Parseval identity for the sign lift of a Boolean function]
\lean{BoolBLR.BLR_soundness_via_fourier_aux_hparseval}
\label{BoolBLR.BLR_soundness_via_fourier_aux_hparseval}
\leanok
\uses{BoolBLR.lift_pm1, BoolFourier.BoolToPM1_sq, BoolFourier.L2_norm_sq, BoolFourier.expectation, BoolFourier.fourier_coeff, BoolFourier.hypercube, BoolFourier.parseval_identity}
\difficulty{2}
Let $f : \{0,1\}^n \to \{0,1\}$ be a Boolean-valued function, and let $g : \{0,1\}^n \to
\bbr$ be its $\pm 1$ lift, given by $g(x) = (-1)^{f(x)}$. Then the Fourier coefficients
$\hat g(S) = \langle g, \chi_S\rangle$, taken over all subsets $S \subseteq [n]$,
satisfy
\[
  \sum_{S \subseteq [n]} \hat g(S)^2 \;=\; 1.
\]
\end{lemma}

\begin{lemma}[Per-coefficient cube bound for functions far from linear]
\lean{BoolBLR.BLR_soundness_via_fourier_aux_h_bound}
\label{BoolBLR.BLR_soundness_via_fourier_aux_h_bound}
\leanok
\uses{BoolBLR.epsilon_far_from_linear, BoolBLR.fourier_coeff_le_of_far_from_linear, BoolBLR.lift_pm1, BoolFourier.fourier_coeff, BoolFourier.hypercube}
\difficulty{2}
Let $f : \{0,1\}^n \to \{0,1\}$ be $\varepsilon$-far from linear, meaning that $0 \le
\varepsilon \le 1$ and every linear function $g : \{0,1\}^n \to \{0,1\}$ (one satisfying
$g(x \oplus y) = g(x) \oplus g(y)$ for all $x,y$) has distance $\Pr_{x}[f(x) \ne g(x)]
\ge \varepsilon$ from $f$. Write $\hat F(S) = \E[F \cdot \chi_S]$ for the Fourier
coefficient of $F = (-1)^{f}$, the $\pm 1$ lift of $f$, at a subset $S \subseteq [n]$.
Then for every such $S$,
\[
  \hat F(S)^3 \;\le\; \hat F(S)^2\,(1 - 2\varepsilon).
\]
\end{lemma}

\begin{lemma}[Linearity gives additivity on XOR]
\lean{BoolBLR.BLR_completeness_aux_h_linear}
\label{BoolBLR.BLR_completeness_aux_h_linear}
\leanok
\uses{BoolBLR.is_linear_bool, BoolFourier.hypercube, BoolFourier.xor_vec}
\difficulty{1}
Let $f : \{0,1\}^n \to \{0,1\}$ be a linear function, meaning that $f(x \oplus y) = f(x)
\oplus f(y)$ for all $x, y \in \{0,1\}^n$, where $\oplus$ denotes componentwise XOR.
Then for all $x, y \in \{0,1\}^n$ we have $f(x \oplus y) = f(x) \oplus f(y)$.
\end{lemma}

\begin{lemma}[Multiplicativity of the sign lift of an XOR-homomorphism]
\lean{BoolBLR.BLR_completeness_aux_h_lift_linear}
\label{BoolBLR.BLR_completeness_aux_h_lift_linear}
\leanok
\uses{BoolBLR.lift_pm1, BoolFourier.BoolToPM1_xor, BoolFourier.hypercube, BoolFourier.xor_vec}
\difficulty{2}
Let $f : \{0,1\}^n \to \{0,1\}$ be a function that respects componentwise XOR, meaning
$f(x \oplus y) = f(x) \oplus f(y)$ for all $x, y \in \{0,1\}^n$, and let $g : \{0,1\}^n
\to \bbr$ be its $\pm 1$ lift, $g(x) = (-1)^{f(x)}$. Then $g$ turns XOR into
multiplication:
\[
  g(x \oplus y) \;=\; g(x)\,g(y) \qquad \text{for all } x, y \in \{0,1\}^n.
\]
\end{lemma}
